\documentclass[11pt]{article}
\usepackage{amsmath,amssymb,amsthm}
\usepackage[margin=1in]{geometry}
\newtheorem{theorem}{Theorem}[section]
\newtheorem{lemma}[theorem]{Lemma}
\newtheorem{proposition}[theorem]{Proposition}
\theoremstyle{definition}
\newtheorem{definition}[theorem]{Definition}

\title{Three facts about the empirical distribution function}
\author{}
\date{}

\begin{document}
\maketitle

\begin{abstract}
We collect an elementary property of the empirical distribution function of a sample, its almost sure uniform
convergence, and the limit law of the Kolmogorov--Smirnov statistic.
\end{abstract}

\section{Setting}

\begin{definition}\label{def:edf}
For $n \ge 1$ and $x = (x_1, \dots, x_n) \in \mathbb{R}^n$, the \emph{empirical distribution function} is
$F_n(t) = \frac1n \#\{ i : x_i \le t \}$ for $t \in \mathbb{R}$.
\end{definition}

When $X_1, X_2, \dots$ are random variables, $F_n$ denotes the empirical distribution function of
$(X_1, \dots, X_n)$, and $D_n = \sup_{t \in \mathbb{R}} \lvert F_n(t) - F(t) \rvert$, where $F$ is the distribution
function of $X_1$.

\section{Results}

\begin{lemma}\label{lem:steps}
For every $x \in \mathbb{R}^n$, $F_n$ is nondecreasing, and $n F_n(t) \in \{0, 1, \dots, n\}$ for every
$t \in \mathbb{R}$.
\end{lemma}

\begin{proof}
If $s \le t$, every $i$ with $x_i \le s$ also satisfies $x_i \le t$, so $F_n(s) \le F_n(t)$. The number
$n F_n(t)$ counts a subset of $\{1, \dots, n\}$.
\end{proof}

\begin{proposition}\label{prop:gc}
If $X_1, X_2, \dots$ are independent and identically distributed, then $D_n \to 0$ almost surely.
\end{proposition}

\begin{proof}
This is the Glivenko--Cantelli theorem.
\end{proof}

\begin{theorem}\label{thm:ks}
If $X_1, X_2, \dots$ are independent and identically distributed with a continuous distribution function $F$, then
$\sqrt{n}\, D_n$ converges in distribution to $\sup_{0 \le u \le 1} \lvert B(u) \rvert$, where $B$ is a Brownian
bridge on $[0, 1]$. The limit has the distribution function
$K(z) = 1 - 2 \sum_{k \ge 1} (-1)^{k-1} e^{-2 k^2 z^2}$ for $z > 0$.
\end{theorem}

\begin{proof}
The convergence is Donsker's invariance principle for the empirical process \cite{donsker}, applied to the
supremum, which is a continuous functional. The form of $K$ is due to Kolmogorov \cite{kolmogorov}.
\end{proof}

\begin{thebibliography}{2}
\bibitem{donsker}
M.~D. Donsker.
\newblock Justification and extension of Doob's heuristic approach to the Kolmogorov--Smirnov theorems.
\newblock \emph{Annals of Mathematical Statistics}, 23:277--281, 1952.
\bibitem{kolmogorov}
A.~N. Kolmogorov.
\newblock Sulla determinazione empirica di una legge di distribuzione.
\newblock \emph{Giornale dell'Istituto Italiano degli Attuari}, 4:83--91, 1933.
\end{thebibliography}

\end{document}
